\documentclass[10pt,letterpaper]{article}
\usepackage[margin=.85in]{geometry}
\usepackage[T1]{fontenc}\usepackage{lmodern,microtype,amsmath,amssymb,booktabs,array,graphicx,xcolor,enumitem,caption,fancyhdr,hyperref,listings}
\definecolor{cernblue}{HTML}{08357E}\definecolor{darkred}{HTML}{7A0101}\definecolor{deepgreen}{HTML}{38761D}\definecolor{navy}{HTML}{17334B}\definecolor{muted}{HTML}{5A6B78}
\hypersetup{colorlinks=true,linkcolor=cernblue,urlcolor=cernblue,pdftitle={Supersonic center of pressure in OpenRocket MIT},pdfauthor={MIT Rocket Team simulation study}}
\setlength{\emergencystretch}{2em}\setlength{\headheight}{14pt}\pagestyle{fancy}\fancyhf{}\lhead{\small\textsc{MIT Rocket Team \enspace|\enspace Simulations}}\rhead{\small OpenRocket MIT 6.6}\cfoot{\thepage}
\captionsetup{font=small,labelfont={bf,color=cernblue}}\setlength{\parindent}{0pt}\setlength{\parskip}{4pt}
\setlist{nosep,topsep=2pt}\setcounter{tocdepth}{1}
\lstset{basicstyle=\ttfamily\small,columns=fullflexible,breaklines=true,frame=single,rulecolor=\color{muted}}
\newcommand{\fig}[3]{\begin{figure}[htbp]\centering\includegraphics[width=#1\linewidth]{figures/#2.pdf}\caption{#3}\end{figure}}
\newcommand{\sub}[1]{\smallskip\noindent{\bfseries\color{navy}#1}\quad}
\makeatletter\let\tabinput\@@input\makeatother
\begin{document}
\begin{center}{\Large\bfseries\color{cernblue} Supersonic center of pressure in OpenRocket MIT}\\[4pt]
{\large Implementation, RASAero II calibration, and the Iris stability discrepancy}\\[6pt]
MIT Rocket Team Simulations \quad$\vert$\quad 9 November 2026\end{center}
\section{Decision and principal result}
The full experimental fin correction places Iris's center of pressure too far aft relative to RASAero II at sustainer burnout. For the preferred 0.50-inch transition, the two codes give the same CG to four decimal places. They put CP at 2.644 and 2.510 m aft of the nose, respectively. That 134 mm difference adds \textbf{1.26 calibers} to OpenRocket's predicted stability margin. The difference before rounding the two margins is 1.255 calibers; the two margins round to 3.91 and 2.65.
 
A trial fitted on five other fin planforms uses 45.9\% of the available force and moment adjustment. It predicts 2.26 calibers at Iris burnout, 0.39 below RASAero, and moves the minimum to within 0.03 s of RASAero. Over the entire sustainer burn the RMS difference in margin falls from 0.95 to 0.44 calibers. Iris did not determine the coefficient.

\textbf{Recommendation: compare with the DATCOM blend, not a global scalar.} A second Java implementation separates fin-loading and forebody corrections using USAF DATCOM tables. It cuts final blind-set CP RMS from 0.58 to 0.44 diameters at zero incidence while keeping two tunnel checks within predefined bounds. On Iris it cuts the burnout overprediction from 1.26 to 0.59 cal (3.24 vs 2.65), though the minimum still comes too early. Bare-body and higher-incidence counterexamples prevent recommending it as the default. The implementation is delivered as a separate research patch; the installed app is unchanged.
The follow-up in Section~\ref{sec:followup} tests a smoother fin handover on Iris, separates three sustainer geometries in both codes, and checks NACA's measured wing--body interference. The new handover improves burn-wide RMS but moves the minimum to 15.53 s. We retain the earlier research blend for general comparison.
Section~\ref{sec:afterbody} tests the airframe itself. Joining equal tubes removes a geometry-dependent CP shift of up to $0.70D$. A DATCOM contracting-boattail model improves the reserved four-fin RASAero comparison at all three incidences. Modeling the stepped center of Iris reduces the burnout prediction to 2.77 cal, while worsening finless-body agreement.
\tableofcontents
\clearpage
\section{Quantities, model definitions and implementation}
All axial stations increase aft from the nose. Let $D$ be the active-stage reference diameter, $S=\pi D^2/4$, $N$ the normal-force slope about zero incidence (or the effective $C_N/\alpha$ at finite incidence), and $J=Nx/D$ its dimensionless first moment about the nose. Then
\[
 X=x_{\rm CP}/D=J/N,\qquad {\rm SM}=(x_{\rm CP}-x_{\rm CG})/D.
\]
A positive CP error at fixed CG overpredicts margin. We report signed errors as OpenRocket minus reference. The moment convention $J$ is a load-position bookkeeping convention; for the tunnel comparison $C_{m_\alpha}=N(X_{\rm ref}-X)$.
 
\sub{Legacy.} The production Barrowman path combines body and fin contributions. At small incidence the supersonic isolated-fin slope is based on $2/\beta$ times exposed one-panel area over $S$, with $\beta=\sqrt{M^2-1}$. The axial load station is computed separately using a Mach-dependent chordwise relation. Counts, instance orientations, interference and mounting station enter the vehicle sum. The legacy simulation option remains the default.
 
\sub{Full finite wing.} Baseline implementation \href{https://github.com/MDNich/ActiveControl_MIT_RktTeam/commit/18bbfbd4}{18bbfbd4} adds a two-face thin-plate slope $4/\beta$ and integrates it over a trapezoid. For span $s$, root chord $c_r$, tip chord $c_t$, sweep distance $h$, span fraction $u$ and chord fraction $v$, the coordinate from the root leading edge is $x=hu+[c_r+(c_t-c_r)u]v$. Set $y=s(1-u)$, $a=x-h$. The downstream tip-relief factor is
\[
 f(x,y)=\begin{cases}1,&a\le\beta y,\\(2/\pi)\arcsin\sqrt{\beta y/a},&a>\beta y.\end{cases}
 \quad
 (N_f,J_{f,\mathrm{local}})=\frac{4}{\beta S}\iint f(x,y)(1,x/D)\,dA.
\]
The midpoint quadrature uses $48\times32$ cells. A per-geometry table interpolates force and first moment every 0.025 Mach from 1.1 to 5. The swept-tip treatment extends the rectangular-wing solution; there is no thickness or fully 3-D shock solution here. See \href{https://farside.ph.utexas.edu/teaching/336L/Fluidhtml/node220.html}{the rectangular-wing derivation} and \href{https://openrocket.sourceforge.net/techdoc.pdf}{OpenRocket technical documentation}, \S3.2.
 
\sub{Where it can contribute.} Define $H(z)=t^2(3-2t)$, $t=\min(1,\max(0,z))$. The handover $w$ is the product of $H$ on the five arguments
\[
 \frac{M-1.1}{0.4},\quad \frac{M s/\sqrt{s^2+h^2}-1}{0.15},\quad
 \frac{\beta-\max(c_t,c_r-h)/s}{0.15},\quad \frac{5-M}{0.35},\quad
 \frac{8^\circ-|\alpha|}{5^\circ}.
\]
Outside $1.1<M<5$, or on non-trapezoidal fins, $w=0$. Legacy and integral loads and their moments are mixed before division. Existing body interference, fin count, azimuth and cant are applied after the single-panel calculation. Body modeling, mass, recovery and drag were unchanged by this trial.
 
\sub{Study implementation.} Let $(N_0,X_0)$, $(N_1,X_1)$ be whole-vehicle outputs with option off and fully on at the \emph{same flight state}. The trial replaces $w$ by $kw$ inside \texttt{FinSetCalc}. Equivalently, because the unchanged terms add linearly,
\[
 N_k=(1-k)N_0+kN_1,\quad
 X_k=\big[(1-k)N_0X_0+kN_1X_1\big]/N_k .
\]
It would be incorrect simply to interpolate the two CP locations. The Java trial class runs ahead of the unchanged JAR on a study-only classpath. A fresh sweep verified both expressions against direct Java outputs at all 39,771 static states to $2\times10^{-12}$ in $D$ and per radian. Integrated flight trajectories are reruns of the solver, not interpolations between two flights.
\clearpage
\section{Experimental design and scoring}
\sub{Independent roles for the data.} The native September atlas contains 92 configurations. We selected 24 with distinct roles and reserved three additional November designs generated in a fresh Windows 11 RASAero II 1.0.2.0 session. We do not treat 92 related designs, or neighboring Mach samples, as independent experiments. The geometry inventory in Appendix A identifies every selected input.
\begin{itemize}
\item \textbf{Fit (5).} Four-inch H-family baseline, reduced sweep, increased span, narrow tip, and unswept. Only zero-incidence samples at $M=1.50,\dots,3.00$ enter training.
\item \textbf{Reserved (7).} Reduced span, intermediate sweep, wider tip, three fins, longer body, and two smaller-diameter S-family rockets. These exercise interpolation and transfer to a different airframe.
\item \textbf{Fresh blind (3).} Square-section fins on the same 4-inch body, with new normalized planforms: short span, wide unswept and highly swept. Geometries and split were specified before their reference exports were examined. The fit was frozen without these cases.
\item \textbf{Controls (5) and sensitivities (7).} Bare bodies isolate residuals with no fin correction. Alternate section, scale and RASAero Rogers-mode variants are sensitivity tests, not additional calibration rockets. The two historical NASA datasets are physical regression checks, previously used to assess the full option; they are independent of the RASAero fit, not newly blinded experiments.
\end{itemize}
\sub{Matched states.} RASAero: open the as-run CDX1, select Aero Plots, File $\to$ Export $\to$ To CSV File. The original exports concatenate angles; select the row's \texttt{Alpha}, then join on Mach, rather than interpolating across angles. Use \texttt{CP} in inches from the nose at that angle, \emph{not} \texttt{CP (0 to 4 deg)}. Multiply by 0.0254 and divide by the CDX1 body diameter in metres. Default CP is \texttt{ModifiedBarrowman=False}; the two files named ``rogers'' deliberately switch it on. See the \href{https://www.rasaero.com/dloads/RASAero%20II%20Users%20Manual.pdf}{RASAero II manual, Aerodynamic Plots}.
 
OpenRocket: construct the same nose, tube and trapezoids in SI units. The fixture supports conical, tangent ogive (parameter 1), von K\'arm\'an (Haack parameter 0) and ellipsoid. Set fin position from the tube bottom using $c_r-\texttt{Location}$; count, span, section, thickness, sweep and chord follow the CDX1. No cant or fillet. For hexagonal blunt-base sections use the MIT triangular leading edge and the CDX1 bevel distance; this is a section approximation. Evaluate both calculators with sea-level ISA, zero rates and zero roll azimuth. Surface drag and motor choice do not enter this static CP call. At $4^\circ$ verify the exported secant against $C_N(4^\circ)/(4\pi/180)$ for these zero-lift-at-zero fixtures.
 
\sub{Fit and audit.} Search $k=0,0.001,\ldots,1$, minimizing the square root of the mean of $(X_k-X_{\rm RAS})^2$ over exactly $5\times151=755$ training states; equal grids assign each rocket equal weight. The grid resolution is numerical precision, not parameter uncertainty. Freeze that $k$ for every remaining comparison. Report bias, mean absolute error, RMS, maximum absolute error and count by geometry, angle and Mach band in the CSVs. Pooled RMS uses state errors, not the arithmetic mean of case RMS values. Audit $0.10$--$0.80$, $0.81$--$1.30$, $1.50$--$3.00$ and $3.01$--$5.00$ separately; $1.31$--$1.49$ is plotted, not scored. Source contains the full grid $0.10$--$5.00$ in 0.01 steps, at $0^\circ,2^\circ,4^\circ$.
 
\sub{Adoption rule.} Require improvement on reserved RASAero CP over the full model, no more than 0.10 $D$ CP RMS deterioration and no more than 5\% slope RMS deterioration on \emph{either} tunnel vehicle. These preselected engineering tolerances are not NASA measurement uncertainties. Fresh holdouts and Iris are reported independently.
\clearpage
\section{Comparison against RASAero}
\fig{.98}{calibration}{The minimum of the training error is interior to the allowed range. The same parameter helps on the reserved group but does not improve the fresh blind group. Each bar uses the same 151 Mach points per rocket.}
\begin{center}\small\begin{tabular}{lrrrr}\toprule
Group & $\alpha$ (deg) & Legacy & Full & Trial $k=0.459$\\\midrule
\tabinput groups.tex \bottomrule\end{tabular}\end{center}
\textbf{Table interpretation.} Numbers above are CP RMS in body diameters for $1.50\leq M\leq3.00$. At zero incidence the trial reduces reserved error from 1.199 to 0.927 $D$ (23\%), but the improvement does not persist at $4^\circ$ or across the three new fin shapes. The five leave-one-training-geometry-out optima are 0.469, 0.473, 0.440, 0.486 and 0.431 in the fit-list order; the narrow range does not imply comparable generalization.
 
The unchanged body-only group is a negative control: all model outputs coincide. Its 0.737 $D$ RMS against RASAero at zero incidence cannot be ascribed to fins. Subsonic and transonic zero-incidence reserved RMS values are 0.13 and 0.91 $D$ for all three models. Above the fitted range ($3.01$--$5.00$), reserved RMS is 2.05 / 1.06 / 0.95 $D$ for legacy / full / trial. These are additional simulator comparisons, not measurements. See \texttt{group-metrics.csv} for every angle and band.
\clearpage
\section{Why a uniform fin multiplier fails}
\fig{.87}{planforms}{Whole-vehicle CP vs Mach at zero incidence. A positive vertical displacement is farther aft. Top: an atlas calibration and two reserved configurations. Bottom: three fresh native exports on previously unseen planforms. Coincident lines are intentional.}
On the blind short-span fin, $\max(c_t,c_r-h)/s=5.44$; $\beta$ reaches only 2.83 by Mach 3. The tip-cone condition never passes. On the highly swept design, $s/\sqrt{s^2+h^2}=0.333$ (approximately); the normal leading-edge Mach only reaches unity just above Mach 3. Both therefore retain legacy predictions throughout the scored band. The short-span result is optimistic by $+2.03\,D$ on average, with $2.18\,D$ RMS; no choice of $k$ can repair it. On the active wide-unswept case, RMS goes from 1.12 (legacy) to 0.24 (full) then back to 0.35 $D$ (trial).
 
\sub{A load diagnostic independent of CG.} At $M=2.50$, $\alpha=4^\circ$, form $N=C_N/\alpha$ with $\alpha$ in radians, and $J=NX$. Subtract the corresponding unfinned H-body values. The increment $\Delta J/\Delta N$ is the axial center of the added load, including installation interference; it is not a direct measurement of an isolated fin.
\begin{center}\small\begin{tabular}{llrrrr}\toprule
Geometry & Source & $N_b$ (/rad) & $X_b$ & $\Delta N$ (/rad) & $\Delta J/\Delta N$\\\midrule
\tabinput body-table.tex \bottomrule\end{tabular}\end{center}
For Hsq, full and RASAero give essentially the same added-load slope and station; scaling the correction down worsens the fin agreement. On the short-span rocket the fin slope is too large; on the wide-unswept rocket it is too small. The body itself also differs. This is why a whole-vehicle CP fit can reward the wrong component change. All eight paired geometries and four Mach stations are in \texttt{body-fin-audit.csv}; the fresh H-body controls have a negligible 0.0001-inch tube-length difference from the archived bare body.
\clearpage
\section{A second model from supersonic aerodynamics literature}
\sub{Why these methods.} Brown's \href{https://ntrs.nasa.gov/citations/19930090940}{NACA Report 839} derives a linearized solution for pointed triangular wings and identifies the leading-edge/Mach-cone ratio as a controlling parameter. It does not supply a general cropped-trapezoid or finite-thickness solution. \href{https://ntrs.nasa.gov/citations/19930082018}{NACA's finite-wing tip analysis} explains why tip relief is important; neither an infinite-wing $4/\beta$ load nor a single constant multiplier can cover all spans. The \href{https://www.pdas.com/datcomrefs.html}{USAF Stability and Control DATCOM} tabulates normal-force slopes for straight tapered wings and potential force and moment for ogive/cone cylinders. These charts offer a practical non-CFD alternative: lookups are cheap enough for an integration step. NASA's \href{https://ntrs.nasa.gov/citations/19780012225}{ODIN module library}, supersonic regime, also separates body potential loading from wing lift. NACA Report 1307, listed in the DATCOM references, addresses the next unresolved term: wing--body interference in lift and pitching moment.

\sub{Isolated wing force.} We extracted the underlying numerical chart entries directly from the \href{https://www.pdas.com/datcomdownload.html}{public Digital DATCOM Fortran distribution}. Mirror the exposed fin about its root to define wing aspect ratio $A=4s/(c_r+c_t)$, taper $\lambda=c_t/c_r$ and $p=\tan\Lambda=h/s$. For positive sweep, Fig.\ 4.1.3.2-56A--F gives ordinate $Q$ at $(\beta/p,Ap,\lambda)$. Convert to panel-area slope as $Q/p$ on the subsonic-leading-edge side and $Q/\beta$ on the supersonic side. For rectangles, Fig.\ 56G supplies $C_{N\alpha}/A$ below $A\beta=1$; above it use $(4-2/(A\beta))/\beta$. For small $Ap$, smooth between the equivalent-rectangle solution at the same $A$ and the chart, attaining the tabulated solution at $Ap=0.25$. This continuation for tapered, nearly unswept fins is ours. The chart-only force uses the legacy fin load station to isolate the slope change. We retain the existing fin/body interference factors. Out-of-table axes saturate at the boundary; reversal of sweep and inverse taper are outside the candidate's supported set.

\sub{Body force and moment.} Let $L_n,d,L_c$ denote nose length, nose base diameter and following cylindrical length. In Figs.\ 4.2.1.1-21A/B, interpolate slope in $(\beta d/L_n,L_c/L_n)$. In Figs.\ 4.2.2.1-23A/B, use the corresponding left or right axis (the reciprocal on the high-$\beta d/L_n$ side); multiply the returned CP fraction by $L_n+L_c$. Limit the represented cylinder to five nose lengths, the chart endpoint; do not multiply a clipped fraction by an unlimited length. Replace the nose potential load and its first moment, retaining the existing incidence-dependent crossflow contributions. Eligibility requires a forward pointed main-stage nose followed by a same-stage, same-radius tube. Cone and tangent-ogive use their charts. Iris's von K\'arm\'an nose uses the equal-length ogive as a \textbf{surrogate}; there is no Haack chart in this port. Power-series noses and ellipsoids retain the baseline body model. Small transitions, tail, pods and other cylinders are not independently replaced. The retained Mach and incidence handovers are engineering choices; the wing charts do not require the old tip-clearance or normal-Mach gates. The port does not implement all of DATCOM, thickness corrections, drag or DATCOM's complete body crossflow model.

\sub{Two coefficients, chosen before the final comparison.} Write a vehicle load pair as $P=(N,J)$. At the same state evaluate $F$ (original full option), $W$ (chart wing with baseline body) and $B$ (chart wing plus chart body). The proposed blend is
\[
           P_* = F+a(W-F)+b(B-W),\qquad X_*=P_{*,J}/P_{*,N}.
\]
Search $a,b=0,0.02,\ldots,1$ on the original five training rockets, discarding pairs outside the two tunnel tolerances of Section 3. The selected pair is \textbf{$a=0.32,\ b=0.40$}; the unconstrained minimum $(1.00,0.44)$ fails the physical constraints. Because the NASA data now constrain the fit, they are no longer independent validation of this blend. These coefficients are empirical weights, not aerodynamic constants. Direct Java vs reconstructed load/moment agreement is better than $2\times10^{-12}$ over 44,190 states.
\clearpage
\section{Transfer results and physical cross-check}
\fig{.98}{transfer-tests}{First five designs train both candidates. The first blind set is fresh for the scalar study, but informed the decision to try DATCOM. Only the final three constitute a new blind test for the literature-based model.}
\begin{center}\small\begin{tabular}{lrrrrrr}\toprule
$M=1.50$--$3.00$, $0^\circ$ & Legacy & Full & Scalar & Wing charts & Both charts & Blend\\\midrule
Training (5) &1.046&.852&.430&.669&.720&.502\\
Reserved (7)&1.557&1.199&.927&1.112&1.226&.857\\
First blind (3)&1.442&1.297&1.305&1.069&.503&.828\\
Final blind (3)&1.259&.579&.487&.426&.986&.438\\
Bare bodies (5)&.737&.737&.737&.737&1.131&.884\\\bottomrule\end{tabular}\end{center}
The blend is more transferable than the scalar on these three held-out groups at zero incidence. Its final-set improvement over full is 24\%. It is \emph{not} the best model on every rocket: at $0^\circ$ near-rectangular fins worsen from 0.35 to 0.62 $D$ while strongly tapered fins improve from 0.88 to 0.34 $D$. At $4^\circ$ the final-group RMS worsens from 0.42 to 0.54 $D$. Wing-only scores best overall on the final $0^\circ$ group, but fails the tunnel slope tolerance. Body-only controls regress at $0^\circ$. See \texttt{chart-cases.csv} and \texttt{chart-groups.csv} for all data, including angles 2 and 4 degrees.

\sub{Physical scale of the tradeoff.} Five supersonic points per vehicle. Column pairs are CP RMS in $D$ and normal-force-slope RMS per radian:
\begin{center}\small\begin{tabular}{lrr|rr}\toprule
&\multicolumn{2}{c|}{ARCAS, $M=1.50$--3.96}&\multicolumn{2}{c}{Yuska B, $M=2.00$--3.47}\\
& CP & $C_{N\alpha}$ & CP & $C_{N\alpha}$\\\midrule
Legacy&2.665&5.671&1.023&4.682\\
Full finite wing&.318&2.882&.126&2.235\\
Scalar trial&1.047&4.373&.511&3.547\\
Both charts, unblended&.859&2.317&.171&2.585\\
DATCOM blend&.312&2.626&.136&2.343\\\bottomrule\end{tabular}\end{center}
The blend meets the preset tolerances. Its reconstructed pitching-moment-slope RMS, which was not a fit constraint, rises from 2.65 to 4.94/rad for ARCAS and from 4.85 to 4.91/rad for Yuska B. A CP fit alone does not guarantee accurate control moments. \texttt{chart-tunnel-metrics.csv} gives both single-method ablations. Sample counts are too small and correlated to justify a flight uncertainty interval.
\clearpage
\section{Iris: the preferred half-inch transition}
Make one change to the clean design used for the historical native export: the first sustainer transition at station 82.00 in is now 0.50 in long. The next transition remains 1.00 in. The updated CDX1 puts the booster at 136.50 in and the total length at 239.25 in. Both engines used digest-matched O6836 and N3366 motor curves. Loaded stack: 73.229153 kg, CG 4.142630 m; loaded sustainer: 26.076044 kg, CG 2.423703 m. The 6.17-inch booster envelope approximates the exposed 6.00-inch motor-case section. Camera shells are equivalent drag areas in RASAero; paired transition diameters are approximations retained from the stable export. Both nozzle exits are unspecified (zero). The original recovery components contribute mass only. Wind is steady at 2 m/s in RASAero; OpenRocket fixes solver and wind seeds separately to 20261108. The first RASAero sustainer sample is 10.53 s, ignition is 15.52 s, and burnout is 20.77 s. The new native export has 16,664 samples.

\fig{.96}{iris-flight}{Preferred Iris: independent time integrations, plotted only after separation. Top: CP and CG measured aft from the nose; the two CG traces nearly coincide at burnout. Bottom: positive static margin in current-stage diameters. Shading marks sustainer thrust.}
\begin{center}\small\begin{tabular}{lrrrrrr}\toprule
&\multicolumn{4}{c}{At sustainer burnout}&\multicolumn{2}{c}{Minimum during burn}\\
Model &$M$&CG (m)&CP (m)&SM (cal)&SM (cal)&Time (s)\\\midrule
\tabinput events.tex \bottomrule\end{tabular}\end{center}
The recommended research blend predicts 3.24 cal at burnout, still \textbf{0.59 cal too high} against RASAero's 2.65. Its burn-wide margin RMS is 0.56 cal, versus 0.95 for full, 0.44 for scalar, 0.34 for the unblended charts, and 1.57 for legacy. The blend also puts the minimum near 18.00 s rather than 20.60 s. A model selected for transfer to other rockets need not be closest on Iris.
\clearpage
\section{Isolating trajectory, geometry and measurement effects}
\fig{.95}{iris-replay}{Static replay on the native RASAero trajectory: each sample sets the same Mach, incidence, and active stage in the OpenRocket aerodynamic API. Positive error means the OpenRocket CP lies farther aft. Axis shows only the sustainer interval.}
The preferred flight comparison uses both clocks and airspeeds as each engine calculates them. For a second check we replay the first 30 s of RASAero states in OpenRocket, reset the active-stage geometry and reference area at separation, set roll azimuth to zero and disable rates. We verify the reference diameters are 6.17 and 4.20 inches; we do not use RASAero's CP to set any replay input. On 526 matched sustainer-burn states the full model has CP bias $+0.53\,D$ and RMS $0.94\,D$. The blend reduces these to $+0.24\,D$ and $0.54\,D$. For the 975 separated states at $M\geq2$ the bias remains $+0.68\,D$ (full: $+1.30\,D$). The residual is aerodynamic/model-form error relative to this native representation; it cannot be removed merely by synchronizing trajectories.
 
CG explains part of the burn-wide margin difference, but not the burnout difference: average OR-minus-RAS CG during thrust is $+0.069\,D$, and at burnout the difference is less than $1\,\mu$m at export precision. In the time comparison we interpolate OpenRocket on the native timestamps within the same stage, never across separation. We check $\delta {\rm SM}=\delta x_{\rm CP}/D-\delta x_{\rm CG}/D$ to $5\times10^{-5}$ caliber. For this half-inch case, model apogees cluster around 25.94 km in OR vs 25.59 km in RASAero; altitude alone would conceal much of the CP error.
 
\sub{Do not mix Iris revisions.} The historical RASAero export used a 0.25-inch transition; the user's saved-flight ORK uses 1.00 inch there. We compared the rocket XML, found that single geometry change and identical simulation conditions, and reran both files as sensitivity cases. For the historical run RASAero gives 2.33 cal at burnout. The corresponding full OR result is about 3.92; the 1.00-inch full OR result also rounds to 3.92. Neither old reference is substituted for the fresh 0.50-inch result. As captured in the native provenance, model translation retains the same simplified taper fields in both CDX1 files: the observed change between whole native runs is not a measurement of a physical transition derivative.
 
\sub{Wind-tunnel provenance.} ARCAS configuration 1 is a half-scale 40.95-in vehicle of 2.25-in diameter. The geometry and aft CP use \href{https://ntrs.nasa.gov/citations/19670020031}{NASA TN D-4014}, Figs.\ 1, 7; experimental series in \href{https://openrocket.sourceforge.net/techdoc.pdf\#page=103}{OpenRocket Fig.\ 6.7} supply the remaining readings. The fixture uses a simplified conical boattail. Yuska B uses the four-fin, 61.7-in, 6-in-diameter model in \href{https://ntrs.nasa.gov/citations/19660004864}{NASA TN D-3182}, Figs.\ 1, 4(b); a three-quarter-power nose and 0.8-degree cant approximate the drawing. The five supersonic readings per vehicle are approximate digitizations, not high-precision coefficient tables. Reference moment stations are 25.95/2.25 and 61.7/6 diameters. Full graph readings, geometry, conventions and omissions are preserved in \texttt{reference/}.
\clearpage
\section{Fin handover, Iris component loads and measured interference}\label{sec:followup}
\subsection*{Can the fin handover recover the time of minimum stability?}
For Iris's sustainer fins, the normal leading-edge Mach number is $M_e=Ms/\sqrt{s^2+h^2}$ with span $s=3.5$ in and sweep $h=5$ in. The prior blend activates the finite-wing pressure solution over $1<M_e<1.15$, or $1.744<M<2.005$. On the native RASAero trajectory its CP error changes from $-0.81\,D$ at 18.00 s to $+0.81\,D$ at 18.60 s. The experiment changes only the fin force and its first moment. It keeps the previously selected forebody weight $b=0.40$.

Let $P_{\mathrm{chart}}$ and $P_{\mathrm{old}}$ be the installed-fin load pairs $(N,J)$ from the chart fin and the prior blend at identical Mach and incidence. The new handover uses
\[
 t=\operatorname{clip}_{[0,1]}\!\left(\frac{M_e-1.15}{0.15}\right),\qquad
 H(t)=6t^5-15t^4+10t^3,\qquad
 P_{\mathrm{new}}=(1-H)P_{\mathrm{chart}}+HP_{\mathrm{old}} .
\]
The chart branch supplies the fin loading through $M=2.005$. By $M=2.267$ the previous prediction is recovered. The width was fixed before scoring. Because $H'(0)=H'(1)=0$, the joining curve matches the force and moment slopes of each branch. We evaluated OpenRocket directly on the half-inch geometry from Mach 1.500 to 3.000 every 0.005 at $0^\circ,2^\circ,4^\circ$, with extra evaluations on either side of both joins. One-sided second-order differences use $\Delta M=10^{-5}$. The check compares the installed load after subtracting the finless result. All 12 force/moment joins met tolerances of $2\times10^{-8}$ in value and $0.01$ per unit Mach in slope. The RASAero reference uses every native 0.01 Mach sample.
\fig{.94}{handover-loads}{Normal force and nose-reference pitching moment of the installed fins at $2^\circ$ incidence. At each Mach subtract the complete finless sustainer from the finned vehicle after applying the same reference area, diameter and moment origin. Shading marks the original leading-edge handover. The dotted vertical line at Mach 2.267 marks the return to the prior blend.}
On 71 fixed Mach points from 1.50 to 2.20 at $2^\circ$, installed-fin $C_N$ RMS error falls from 0.0292 to 0.0107. Nose-reference $C_m$ RMS error falls from 0.923 to 0.346. At $4^\circ$, the corresponding changes are 0.0506 to 0.0099 and 1.597 to 0.327. During the matched-state replay, the new CP errors at 18.00 and 18.60 s are $+0.30$ and $+0.34\,D$. Across the 59 powered native states with $1<M_e<1.15$, CP RMS decreases from 0.587 to 0.309 $D$. Across the entire burn it decreases from 0.541 to 0.451 $D$. At 20.77 s both predictions give $+0.57\,D$ on the same RASAero state. The two independent OpenRocket flight integrations differ slightly in burnout Mach, so their burnout margins differ by 0.0023 cal.
\clearpage
\subsection*{Iris flight and regression checks}
\fig{.95}{handover-flight}{Static margin during sustainer thrust. Dots mark the minimum over the entire burn, including the ignition transient. RASAero reaches its minimum at 20.60 s, shortly before burnout.}
\begin{center}\small\begin{tabular}{lrrrrr}\toprule
Model & Minimum (cal) & Time (s) & Burnout (cal) & Burn RMS (cal) & Bias (cal)\\\midrule
RASAero II &2.581&20.600&2.650&\multicolumn{2}{c}{reference}\\
Prior blend &2.767&18.000&3.240&.563&+.168\\
New handover&2.990&15.530&3.243&.439&+.273\\
Unblended fin and body charts&2.201&20.603&2.226&.344&$-$.260\\\bottomrule\end{tabular}\end{center}
The new handover reduces RMS margin error by 22\%, but its minimum occurs 5.07 s early, near ignition. Accuracy at minimum margin, signed bias, and accuracy over the whole burn favor different choices. The local crossover contributes to the shape of the old prediction; smoothing it alone does not recover the timing or eliminate the high-Mach residual.

For the existing seven reserved rockets at zero incidence and Mach 1.50--3.00 the CP RMS changes from 0.857 to 0.911 $D$. On the original three blind rockets it changes from 0.828 to 0.837 $D$, and on the final three from 0.438 to 0.434 $D$. Those geometries are regression cases in this pass. At the original ARCAS and Yuska B supersonic sample points the two models coincide to the precision of the stored output. Based on the flight timing, reserved geometries and interference results below, we retain the new handover as an experimental branch. The published research patch continues to represent the earlier blend.
\clearpage
\subsection*{Where the Iris burnout discrepancy originates}
The linked simulation study ran five native static exports in RASAero II 1.0.2.0. Three correspond to the bare 32-in Haack nose and first 50-in tube, the entire 136.5-in finless sustainer, and that sustainer with its four fins. The two additional runs remove the equivalent camera protuberance from the finless and finned geometries. That change gives identical native $C_N$, potential $C_N$, and CP at all 906 matched states across both controls. In OpenRocket we repeated the experiment both on the actual half-inch ORK and on a fixture built from the literal CDX1 external dimensions. The ORK contains explicit camera pods and tapered steps; the CDX1 represents the camera by area and the two short transitions by equal diameter fields. Both sets of OpenRocket results are retained.

Use the nose tip as the moment origin and $D=4.20$ in throughout. For the bare forebody, convert the native normal coefficient to the common area by multiplying by $(4.02/4.20)^2$. At positive incidence use the native angle-specific $N=C_N/\alpha$, $J=N x_{\rm CP}/D$. Compute successive differences in $(N,J)$ across the three fixtures. The installed-fin increment includes interaction. Each increment's exact contribution to whole-vehicle CP error is $(\Delta J_i-X_{\rm RAS}\Delta N_i)/N_{\rm OR}$, where the deltas are OR minus native RASAero, $X_{\rm RAS}$ is the finned vehicle's CP in diameters and $N_{\rm OR}$ its OpenRocket slope. Summing these terms recovers the total CP difference.

\begin{center}\small\begin{tabular}{lrrrrr}\toprule
At $M=2.80$, $2^\circ$ & $N_{\rm RAS}$ & $N_{\rm OR}$ & $J_{\rm RAS}$ & $J_{\rm OR}$ & CP contribution ($D$)\\\midrule
Bare nose and tube &2.731&2.876&15.55&17.50&$-$.146\\
Additional finless airframe &1.536&.809&33.65&20.12&+.331\\
Installed fins &5.359&5.489&170.21&174.39&+.134\\\midrule
Whole vehicle &9.627&9.174&219.41&212.01&+.319\\\bottomrule\end{tabular}\end{center}
The largest positive term comes from adding the rest of the airframe. It includes the aft tubes, steps and camera representation. Removing the explicit OpenRocket camera pods in an extra matched pair changes the contribution by $-0.014\,D$ at $2^\circ$ (pods present minus absent). The aft tubes and steps therefore account for approximately $+0.345\,D$ of the error at this state. On the CDX1-dimension OpenRocket fixture the three terms are $-0.149,+0.377,+0.136\,D$, respectively. At $4^\circ$, the actual ORK gives $-0.995,+0.485,-0.097\,D$. The forebody then dominates the signed error.

\fig{.91}{handover-components}{Native zero-incidence CP at each Mach. All three panels use the same 4.20-inch unit. At Mach 2.80 the native finless position is $9.09\,D$ aft of the nose versus $5.49\,D$ in the original OpenRocket geometry. Bare forebody CP differs by only $0.18\,D$.}
\clearpage
\subsection*{Load attribution at the flight's near-zero incidence}
At the 20.77-s native burnout sample, $M=2.8081$ and $\alpha=0.000434^\circ$. The static CSV publishes $C_N=0$ at exactly zero incidence. Its column labeled 0--4-degree slope is a finite secant. To estimate the derivative relevant to this flight, pair the native zero-angle CP with the potential-force secant estimated at $2^\circ$ for each fixture. On 330 matched low-incidence supersonic burn samples the fully assembled proxy has a slope RMS difference of $0.105$/rad from the native flight value $(C_L\cos\alpha+C_D\sin\alpha)/\alpha$. At burnout this comparison gives 8.397 versus 8.416/rad.

At Mach 2.80 on the actual ORK the resulting three contributions are $+0.036$ (forebody), $+0.391$ (additional airframe), and $+0.146\,D$ (installed fins). Using the potential value at $4^\circ$ reproduces all three to the shown precision. As a stress test, substituting the \emph{total} $2^\circ$ secant yields $+1.67,+1.04,-2.13\,D$ with the same total. The choice of force derivative matters for individual terms even though the zero-incidence CP positions themselves are directly exported.
RASAero's static test interface exported 1.00 degree when given 0.01 degree. We exclude that setting from the near-zero comparison. The finite-incidence table above gives direct native results; individual zero-incidence shares in this section depend on the stated slope estimate.


\clearpage\subsection*{Physical test definition}
We transcribed the experimental entries of Tables I and II in \href{https://ntrs.nasa.gov/citations/19930084523}{Nielsen, Katzen and Tang, NACA TN 3795 (1956)}. The six complete triangular wings have aspect ratios 0.67, 1.34, 2.02, 2.69, 3.33 and 4.00. The split was recorded before inspecting Table II. W1/W3/W5 form the development set and W2/W4/W6 the reserved set, each at Mach 1.50 and 2.02. The candidate used the same coefficients on every test. Repetition of one shared body reading across wing sizes does not increase the sample count.

The body has length 11 in, diameter 1.50 in and a 15-degree semiangle cone tangent to a circular fairing of radius 11.02 in. We represented the fairing as 24 successive frusta. All trailing edges lie 8 in aft of the nose. For the isolated wing the full triangle extends to the centerline on a negligible-radius carrier. For the assembled configuration each exposed semispan is $s_m-0.75$ in with a 3-in root chord. We evaluated two opposing fins in their lift plane and matched the tunnel Reynolds number of $5.5\times10^6$ based on body length. OpenRocket approximates the tested double-wedge section of 8\% local thickness with a constant-thickness rounded plate. This affects the axial-drag correction and limits the precision of the isolated-wing comparison. Doubling the fairing resolution and halving the carrier radius changed the predicted load station by at most 0.00053 in and the base-area normal-force slope by less than 0.00004/rad across all fixtures.

Convert measured lift slope from per degree to per radian and compute $C_{m\alpha}=(dC_m/dC_L)C_{L\alpha}$. Use table area $S$ and mean chord $\bar c$ for wing and combination, the extended triangle centroid $x_r=8-c_r/3$ as their moment origin, and the body base area, length and nose for the body alone. With OpenRocket's common base-area slope $N$ and zero-incidence axial coefficient $C_{A0}$, the predictions on the published references are
\[
 C_{L\alpha}=(N-C_{A0})S_b/S,\qquad
 C_{m\alpha}=N(S_b/S)(x_r-x_{\rm CP})/\bar c .
\]
We compare these directly to Table II before interpreting $x_r-\bar c C_{m\alpha}/C_{L\alpha}$ as an effective lift-center position.
\fig{.88}{handover-tn3795}{Measured and predicted assembled-vehicle lift and centroid-reference pitching-moment slopes. Open and filled markers identify the preassigned development and reserved geometries. Each reserved score below uses six measurements over three geometries and two Mach numbers.}
\clearpage\subsection*{Independent force, moment and interference results}
\fig{.89}{handover-interference}{Subtract the isolated body and complete wing from the assembled measurements on common references. The observed lift increment approaches zero as span increases. Both OpenRocket variants give a larger loss and a larger nose-reference moment change over much of the tested range.}
\begin{center}\small\begin{tabular}{lrrrr}\toprule
&\multicolumn{2}{c}{Lift-slope RMS (/rad)}&\multicolumn{2}{c}{Moment-slope RMS (/rad)}\\
Reserved cases & Prior & Handover & Prior & Handover\\\midrule
Isolated full wing &.307&.229&.117&.105\\
Assembled wing and body &.796&.555&.122&.109\\
Interference increment &.435&.547&1.385&1.602\\\bottomrule\end{tabular}\end{center}
The installed configuration and the isolated wing improve on both native observables. The assembled effective lift-center RMS improves from 0.330 to 0.255 body diameters. For the shared body alone the two models coincide. Measured lift slopes at Mach 1.50/2.02 are 1.95/2.64 per radian, while the predicted values are 1.56/1.65. The corresponding measured nose-reference moment slopes are $-0.390/-0.527$, versus $-0.379/-0.379$.

For the interference row subtract body-alone and full-wing-alone dimensional loads from the combined loads, then normalize by the complete wing $S$ and $S\bar c$. Shift to a common lift-based nose datum using $S(\bar c\,C_{m\alpha}-x_r C_{L\alpha})$ for model and measurement. The derived row therefore uses the paper's lift-centroid convention; exact body-axis translation would also require each measured zero-incidence axial force. On the reserved wings the mean lift interference is $-0.264$/rad in the measurements and $-0.789$/rad in the new model. The corresponding mean common-datum moment increments are $+0.113$ and $+1.661$/rad. The model overstates the loss of lift and the change in moment when the delta wing is mounted. Better total CP can coexist with larger errors in the interaction itself.

The frozen split and all per-configuration outcomes, including development failures, are in \path{followup/protocol.md} and \path{results/followup/}. The top-level reproduction driver reruns the Java flight, fixed-incidence sweeps, ablations, tunnel tests and figures; \path{inputs/tn3795/source-extract.pdf} preserves the source tables and drawing.

\clearpage\section{Afterbody loads, body segmentation and the Iris sustainer}
\label{sec:afterbody}
A body tube can occupy the same space whether a designer draws it as one piece or splits it at a coupler. In the previous research blend, the nose chart read just the first tube length. On a 12-inch ogive with a 48-inch, 4-inch-diameter cylinder, splitting the cylinder into four pieces changed the predicted CP by up to $0.696D$ at zero incidence. Joining equal-diameter, touching tubes for the chart lookup reduced the maximum difference below $10^{-8}D$ over all 453 tested Mach/incidence combinations. This correction addresses the representation of the same exterior in the model.

\sub{A second hypothesis concerns the taper at the base.} A contracting boattail supports a negative normal-force increment. Even a small error in this load matters when it acts many diameters behind the nose. On a bare body the positive forebody and negative tail loads can nearly cancel. CP can then lie forward of the nose and be especially sensitive to errors in either load. For this reason we test the force and first moment of the taper against a straight control as well as the CP of the whole rocket.

\sub{Methods tested.} The prior blend fixes the fin and nose chart weights at 0.32 and 0.40. We evaluated three additional methods individually and jointly. The code extracts the original Fortran DATA arrays from the archived \href{https://www.pdas.com/datcomdownload.html}{USAF Stability and Control Digital DATCOM distribution}, subroutine \texttt{SUPBOD}. Every new table uses multilinear interpolation with the first Fortran index contiguous. No coefficient was fitted to the new rockets.

\begin{enumerate}[leftmargin=*]
\item \textbf{Center section.} For an eligible nose of length $L_n$ and base diameter $D_n$, group the continuous, same-stage center section up to the last eligible tube. The extended model permits monotone conical steps whose diameters remain within 10\% of $D_n$. With end diameter $D_1$, set $\bar D=(D_n+D_1)/2$, $u=\sqrt{M^2-1}\,D_n/L_n$, and $f=\min(5,L_cD_n/(L_n\bar D))$. Read normal slope and CP fraction from figures 4.2.1.1-21A/B and 4.2.2.1-23A/B. Normalize slope on the area of $D_1$ and multiply CP fraction by $L_n(1+f\bar D/D_n)$. Blend normal force and its nose-referenced first moment. Retain the unblended fraction of the incorporated step loads. Applying this smooth-body construction across short stepped transitions is an engineering approximation. The equal-tube variant stops at the first change in radius or non-tube component.
\item \textbf{Terminal taper.} For taper length $L$, fore/aft diameters $d_1,d_2$, $q=d_2/d_1$, $\beta=\sqrt{M^2-1}$ and taper half-angle $\theta$, use the fore-area coefficient. For $q<1$, figure 4.2.1.1-22A gives $\Delta N=(1-q^2)F(L/(\beta d_1))$. Figure 4.2.2.1-24 gives its station from the start of the taper, $x_t=L H(\beta\tan\theta,q)$. For a flare, figure 4.2.1.1-22B gives $\Delta N=G(q)\cos^2\theta$. The tabulated moment $K$ from figure 4.2.2.1-25A gives $x_t=-K(\theta,1/q)d_1q^3/\Delta N$. Convert area, shift the moment to the rocket datum and blend force and moment. The boattail port uses the chart's local load station and translates it by the taper start in OpenRocket. This is a component implementation of the charts.
\item \textbf{Crossflow.} Replace the constant 1.1 in the symmetric body's finite-incidence lift term by the crossflow coefficient in figure 4.2.1.2-35B evaluated at $M|\sin\alpha|$. Keep OpenRocket's projected area, centroid and $\sin^2\alpha$ term. This correction vanishes at zero incidence.
\end{enumerate}

\sub{Eligibility.} Apply the established smooth envelope over Mach 1.1--1.5 and 4.65--5.0, with full incidence weight through $3^\circ$ and zero at $8^\circ$. The taper must be the last active centerline part and touch a radius-matched tube. Use chart limits $\beta\tan\theta\leq0.5$ and $L/(\beta d_1)\leq2.5$ on contractions. Smoothstep to zero over the final tenth of each interval. An abrupt switch at the first limit was visible on the ARCAS Mach sweep before this correction. Flare queries require $1<q\leq3$ and $5^\circ\leq\theta\leq90^\circ$. The stage mask invalidates cached component eligibility when a booster separates.

\clearpage
\sub{Native RASAero campaign.} RASAero II 1.0.2.0 in the Windows 11 VM produced independent Aero Plots exports for 13 configurations. Three existing native Iris component exports complete the comparison. Each row below uses 151 Mach values from 1.50 to 3.00 in 0.01 increments at each incidence. The reference is the native angle-specific CP in inches aft of the nose, normalized by that configuration's maximum body diameter. P denotes the prior blend. N joins equal tubes and applies DATCOM only to terminal contractions. A includes the stepped-center, terminal taper and crossflow methods together.
\begin{center}\small\setlength{\tabcolsep}{4.5pt}\begin{tabular}{l rrr rrr rrr}\toprule
&\multicolumn{3}{c}{$0^\circ$}&\multicolumn{3}{c}{$2^\circ$}&\multicolumn{3}{c}{$4^\circ$}\\
Case & P & N & A & P & N & A & P & N & A\\\midrule
\tabinput afterbody-static.tex \bottomrule\end{tabular}\end{center}

\sub{Geometry controls.} A has a 12-inch tangent ogive and 48-inch afterbody of nominal diameter 4 inches. Its 4-inch terminal segments end at 3.0 or 4.8 inches. Its internal-step alternative has a 36-inch tube, 0.50-inch step to 4.2 inches, then an 11.5-inch tube. B uses a 9-inch conical nose and 36-inch, 3-inch-diameter afterbody. Its 3-inch tapers end at 2.4 or 3.6 inches. The three-fin B pair shares the same 4.5/2-inch root/tip, 2.25-inch span, 2.5-inch sweep, 0.125-inch square section and fore-aft placement. Family B was reserved before the A/B calculations. After inspecting those results we fixed the narrower method and requested a fresh confirmation pair C. C has a 5-inch diameter, 18.5-inch tangent ogive, and 49-inch afterbody. In the variant the final 6 inches taper to 4.5 inches. Four fins of 6/2-inch root/tip chord, 3.5-inch span, 4.5-inch sweep and 0.15-inch thickness occupy the same root stations, 54--60 inches, on both models.

On the straight A and B controls, the crossflow coefficient is the only active change in the all-method model. For A it improves CP RMS from 0.33 to $0.24D$ at $2^\circ$ and worsens it from 0.38 to $0.50D$ at $4^\circ$. B shows the same reversal. This incidence dependence argues for retaining the original crossflow law in the narrow variant.

The straight Iris controls remove the steps while retaining the 136.5-inch length, 4.02-inch cylinder and original fin root stations at 128.5--136.5 inches. The static reconstruction uses the preceding component's aft radius as the front of a native transition and \texttt{Transition/Diameter} as its outlet. In two additional VM exports, changing only \texttt{RearDiameter} left the coefficients identical. Changing only \texttt{Diameter} moved the rear OD in the editor and altered all 900 checked CP values. A separate \texttt{BoatTail} stores its outlet in \texttt{RearDiameter}. Equal values in the two serialized diameter fields can describe a real taper. This resolves an ambiguity in the earlier literal-CDX1 diagnostic. Each terminal \texttt{BoatTail} contributes geometry once. The body-tube boattail fields duplicate its dimensions. The OpenRocket static fixtures reproduce the external axisymmetric profile and attached fins. The preceding native camera-area removal controls gave identical normal loads and CP at all 906 matched states. The integrated flight uses the actual ORK with its pods and motors.

\clearpage
\sub{Force and moment of a geometry change.} For each matched pair use the same $D_*$ and area $S_*$ for both configurations. At finite incidence recover OpenRocket $C_N=\alpha N$ with $\alpha$ in radians. For each code separately compute
\[
 \Delta C_N^*=\left[C_N(D/D_*)^2\right]_{\rm variant-control},\qquad
 \Delta J^*=\left[C_N(D/D_*)^2x_{\rm CP}/D_*\right]_{\rm variant-control}.
\]
The table gives RMS error in these increments at $2^\circ$. Use $D_*=4$, 3, 5 and 4.20 inches for A, B, C and Iris. These nose-referenced first moments use aft-positive station. The same analysis at $4^\circ$ and all component-only ablations are in the result CSVs.
\begin{center}\small\begin{tabular}{lrrr rrr}\toprule
&\multicolumn{3}{c}{RMS error in $\Delta C_N^*$}&\multicolumn{3}{c}{RMS error in $\Delta J^*$}\\
Variant & P & N & A & P & N & A\\\midrule
\tabinput afterbody-pairs.tex \bottomrule\end{tabular}\end{center}

\sub{Wind-tunnel evidence.} ARCAS contains a contracting afterbody. Across its five supersonic measurements the candidate changes CP RMS from 0.312 to $0.363D$, normal-force-slope RMS from 2.63 to 2.31/rad, and pitching-moment-slope RMS from 4.94 to 3.08/rad about the benchmark datum. Yuska B has a straight afterbody and gives the same $0.136D$ CP RMS. These measurements give a mixed assessment of boattail transfer.

For a separate flare check we used Jorgensen, Spahr and Hill's \href{https://ntrs.nasa.gov/citations/19720065133}{NASA TM X-653}, figures 1(b) and 4(a,b). Both configurations have a 1.25-inch-diameter midbody, blunted $14^\circ02'$ cone, a $2D$ cylinder and a $1.37D$ terminal segment. One terminal segment is straight and the other flares to $1.48D$. The numerical forebody uses a sphere of radius $0.25D$ tangent to the cone, with 40 frusta over the spherical portion. We read the plotted small-angle slopes and CP at approximately Mach 1.97, 2.94 and 4.06. Normalize both force slopes to the midbody area, convert degrees to radians, and measure CP and $J$ from the nose--cylinder junction. The flare's total CP RMS improves from about 0.074 to $0.053D$. The measured flare-minus-straight increment gives the opposite assessment. Its slope RMS rises from 1.06 to 1.10/rad and its first-moment RMS rises from 2.26 to 2.57/rad. The straight control already has $0.95D$ CP RMS error. The original plot and digitization are archived. Two decimal places describe graph readings only.

\sub{Checks on the implementation.} Each fixed-incidence case uses a single solver across its Mach grid. A detach/reattach test compares a reused calculator with a fresh one at each stage state. The earlier ARCAS and Yuska B fixtures were scanned every 0.001 Mach from 0.750 to 5.001, checking finite loads and adjacent changes in CP and slope. The packaged narrow JAR reproduces 8,613 harness rows with a maximum CP difference below $2\times10^{-12}$ inches compared with the class-directory run. Legacy flight histories for both research runs are byte-identical to the previously archived legacy flight. The prior-mode Iris replay reproduces the earlier shared-state solver output to $10^{-9}$. The segmented cylinder's one-, two- and four-tube representations also agree in normal slope to $10^{-8}$/rad with the geometry correction.

\clearpage
\sub{Iris in two complete flight solvers.} The preferred first transition is 0.50 inches long. We reran the full two-stage flight in the Windows VM with the original O6836 and N3366 motors. All 16,664 rows and all 24 exported columns reproduced the archived trajectory byte for byte. The OpenRocket results below are separately integrated flights using identical seeds. Within the 15.52--20.77 s sustainer burn we interpolated the separated-stage OpenRocket histories at the 526 native timestamps. Minimum and its time use the integration samples including interpolated endpoints. RMS and bias are in calibers and take OpenRocket minus RASAero.
\begin{center}\small\begin{tabular}{lrrrrr}\toprule
Model & Burnout & Minimum & Time of min (s) & Burn RMS & Burn bias\\\midrule
\tabinput afterbody-flight.tex \bottomrule\end{tabular}\end{center}
\fig{.98}{afterbody-comparison}{Left, zero-incidence CP error against native Aero Plots. Right, integrated Iris stability through the sustainer burn. The narrower variant overlaps the prior blend to plotting precision on this flight.}

\sub{Attribution on Iris.} Replaying the same native Mach, incidence and stage at burnout reduces the CP error from $+0.572D$ to $+0.112D$. The stepped-center method alone gives $+0.112D$, the taper alone $+0.572D$ and the crossflow alone $+0.572D$. The complete sustainer ends in a tube. On the 975 sampled sustainer states above Mach 2 the all-method CP RMS drops from $0.680D$ to $0.285D$. Yet on the stepped finless sustainer the zero-incidence static CP RMS rises from $2.86D$ to $4.16D$. At $2^\circ$ the matched Iris step's first-moment discrepancy also grows. The lower assembled burnout margin therefore cannot establish the assumed pressure distribution along the steps. The minimum remains about 2.6 s early and is now 0.23 cal below RASAero.

\sub{Use in MIT 6.6.} Aggregate geometrically identical tubes. Offer contracting boattails as an experimental comparison for configurations within the chart envelope, with the ARCAS CP regression visible. Keep the stepped-center method available to bound model sensitivity on Iris. Keep the previous flare and crossflow terms for general use. The separate patch and JAR in \texttt{implementation/} select the narrow method by default when the experimental fin option is active. In a fresh JVM use \texttt{-Dstudy.afterbody=} followed by \texttt{previous}, \texttt{segments}, \texttt{center}, \texttt{tail}, \texttt{crossflow}, \texttt{combined}, or \texttt{narrow}. All modes preserve the published fin-blend constants. The tables are component ports of selected DATCOM methods, with a geometry extension expressly identified above. The comparisons cover static longitudinal loading at small incidence.

\clearpage\section{Engineering recommendation and reproducibility}
Use the prior blend as a comparator, aggregate equal-diameter tubes, and evaluate the terminal boattail variant for in-domain contractions. The ARCAS CP regression limits confidence in applying it generally. For the preferred Iris compare the new native 2.65-cal burnout with the prior blend's 3.24 cal and the stepped-center experiment's 2.77 cal. The broad experiment improves agreement at burnout and underestimates the minimum. Its poorer finless-body and step-moment results constrain interpretation of the assembled rocket. The TN 3795 tests in Section~\ref{sec:followup} also identify an interference-moment error. High incidence, spin, multibody shock interactions and unusual sections remain outside these checks.

\sub{Code map.} \texttt{FiniteWingPressure} integrates panels; \texttt{FinSetCalc} applies instance loading; the Barrowman calculator selects component calculators. The original option, XML load/save and Swing checkbox are in commit \texttt{18bbfbd4}. The study patch adds \texttt{ChartData}/\texttt{ChartModel}, blends the two fin load pairs in \texttt{FinSetCalc}, and replaces only the eligible nose potential part in \texttt{SymmetricComponentCalc}. The existing experimental flag activates the study variant in a patched process; flag-off remains legacy. No change is made to the installed application. The study source includes a patch and a Java overlay JAR; \texttt{git apply --check} passed against the pinned source. All reported variant flights run actual OpenRocket integration code.

\sub{Run order.} From the \href{https://github.com/MDNich/MIT-RT-26_27-Work/tree/main/design_work/openrocket_CP}{complete team study directory}, use JDK 17, Python with NumPy and Matplotlib, and a TeX installation with pdfLaTeX, PGF, Beamer and Latin Modern. The run here used Java 17.0.21, Python 3.15.2, NumPy 2.5.4 and Matplotlib 3.11.2. The reference JAR SHA-256 and raw-input hashes are in the README, CSV manifest and native export notes.
\begin{lstlisting}
git clone https://github.com/MDNich/ActiveControl_MIT_RktTeam.git OR
git -C OR checkout 18bbfbd4
(cd OR/clone/openrocket && ./gradlew shadowJar)
export OR_REPO="$(cd OR && pwd)"
cd /path/to/openrocket_CP
bash reproduce.sh
\end{lstlisting}
The driver reconstructs the archived atlas subset, compiles the Java sweeps, runs paired Iris flights, fits the scalar, checks it against direct Java, builds wing/body ablations, fits the constrained grid, scores the locked final set, replays the preferred trajectory, then regenerates all tables and figures and both PDFs. It does not rerun Windows software. To regenerate native references, use the CDX1 and motor files with RASAero II 1.0.2.0 and the exact GUI procedures and settings in each native folder. No proprietary binary is bundled. The exported native data suffice for independent scoring. The report-source and slide-source archives also rebuild their PDFs directly without rerunning a solver.

\sub{What to inspect first.} \texttt{design.json} fixes the first split and scoring bands. \texttt{literature/protocol.md} records the second pass and final holdouts. \texttt{results/transfer-grid.csv} contains all 2,601 parameter candidates. \texttt{chart-groups.csv}, \texttt{chart-cases.csv} and \texttt{chart-tunnel-metrics.csv} give all scored outcomes. \texttt{iris-events.csv}, \texttt{iris-errors.csv} and \texttt{iris-common-state.csv} separate flight, sampling and aero effects. All plots use Matplotlib PGF with pdfLaTeX typesetting, using the supplied blue \#08357E, dark red \#7A0101 and green \#38761D palette.
\clearpage
\appendix\section{Geometry ledger and per-case results}
Dimensions in inches. $D$ is body diameter, $L_n/L_b$ lists nose and tube length, $n_f$ is fin count, $c_r/c_t$ are root and tip chords, and $s/h$ are exposed span and axial sweep. The last three columns give $0^\circ$, $M=1.50$--3.00 CP RMS against RASAero. The two physical benchmark geometries and Iris are separate. Geometry is taken from each as-run CDX1, not inferred from the case name.
\begin{center}\scriptsize\setlength{\tabcolsep}{3.5pt}\begin{tabular}{llrrrrrrrr}\toprule
Case&Role&$D$&$L_n/L_b$&$n_f$&$c_r/c_t$&$s/h$&Legacy&Full&Scalar\\\midrule
\tabinput matrix.tex \bottomrule\end{tabular}\end{center}
\sub{Final confirmation designs.} Each has an ogive nose three diameters long, a tube twelve diameters long and a 0.5-inch trailing-edge setback. Hashes of the model files were frozen before native export. These last designs were not used to choose the coefficients.
\begin{center}\small\begin{tabular}{lrrrrrr}\toprule
Case & $D$ & $n_f$ & $c_r/c_t$ & $s/h$ & Full RMS & Blend RMS\\\midrule
Moderate sweep &3.25&3&5.25/2.375&3.125/2.25&.33&.28\\
Strong taper &4&4&7.50/.75&4.125/3.75&.88&.34\\
Near rectangle &4&4&4.75/4.50&3.875/.125&.35&.62\\\bottomrule\end{tabular}\end{center}
The native model files also retain thicknesses, nose shape, finish, bevels and locations. Half/double scale and alternate CP modes furnish sensitivity checks but do not count as independent geometry wins. The two tip-length variants also change sweep relative to Hsq; they are not clean one-variable taper experiments.
\end{document}
